\section{Discussion and open problems}

Order $10$ is excluded by the complete master searches in
Computational Theorem~\ref{cthm:order10-nonexistence}. The explicit
order-$12$ witness refutes the former power-of-two existence conjecture.
Neither fact decides order $18$. The earlier subset-action, sign,
incidence, and flag-surface investigations remain structural results
and diagnostic boundaries; their unresolved solver runs are not
additional proofs of nonexistence.

The most concrete recent advance is constructive rather than an
order-$18$ obstruction. Theorems~\ref{thm:row-fibred-pattern}
and~\ref{thm:affine-fibre-orbits} give an exact pattern formula and
a within-family equivalence criterion. Their numerical consequences
distinguish many main classes, but count only a specified family.
They do not imply that every FFF square has a recoverable binary
quotient or that arbitrary direct products cancel.

The position--value channel of Theorem~\ref{thm:value-color-channel}
reduces the formula size without removing any view. The finite
order-$18$ exclusions in Section~\ref{sec:order18-scope} instead close
specified construction routes. These are different forms of
progress: an equivalent model is not a solved search, and exclusion
of a route is not a complete order classification.

\begin{enumerate}
\item Decide existence at order $18$, without replacing unrestricted
Latin squares by an unjustified affine, quotient, symmetry, or trade ansatz.
\item Determine minimal dyadic exponents $h(m)$ beyond the cases
already settled, and identify which upper bounds can be sharpened.
\item Find a structural explanation of the order-$10$ exclusion
that does not require its exhaustive master graphs.
\item Determine which constructions produce classes outside the
fixed affine and row-fibred families, and when weaker hypotheses
still recover their quotient layers.
\item Complete a systematic priority comparison and independent
expert review, including the computer-assisted dependencies.
\end{enumerate}

These questions concern the even-cycle condition in three line views.
They are not a resolution of the global Alon--Tarsi
conjecture~\cite{AlonTarsi1992}. The research archive retains the
historical investigation sequence; this edition separates that history
from the currently supported theorem statements.
